\scanpage{65}

\margin{\fbox{$\unsure{\delta|N}\overset{\mathrm{dfn}}{=}\operatorname{Im}(\operatorname{Pic}(X_0)\longrightarrow\underline{\operatorname{NS}}_{X_0/k}(k))$}}
\uline{Théorème}\quad $S$ trait str. local

$X$ schéma propre sur $S$, régulier, tel que $\underline O_S\xrightarrow{\sim}f_*(\underline O_X)$, et que $X_\eta$ admette \inserted{un 0-cycle de degré 1}

$\Delta$ le groupe des diviseurs sur $X$ tels que $\operatorname{supp}\unsure{\Delta|D}\subset X_0$
\[
c:\Delta\mathbin{-}\deleted{\unsure{\operatorname{NS}(X_0)}}\ \underline{\operatorname{NS}}\deleted{\unsure{{}_{X_0/k}(k)}}(X_0)
\]
défini comme composé
\diagram{\begin{tikzcd}[column sep=small]
\Delta\arrow[r]&\operatorname{Pic}(X)\arrow[r,dash]&\operatorname{Pic}(X_0)\arrow[r,dash]&\underline{\operatorname{Pic}}_{X_0/k}(k)\arrow[r,dash]&\underline{\operatorname{NS}}_{X_0/k}(k)\arrow[r]&N
\end{tikzcd}}
\[
\left\{\begin{gathered}
\delta=\deleted{\unsure{\operatorname{Coker}c\text{ qui}}}\ \text{ss-groupe de torsion de }\operatorname{Coker}c,\\
\deleted{\unsure{\operatorname{NS}\text{ et }Z}\ \text{\unsure{qui est le}}\ \text{\illegible{2}}\ \operatorname{NS}(X_0),\ \text{de sorte que }\Delta\longrightarrow N\xrightarrow{\lambda}\delta\longrightarrow0}
\end{gathered}\right.
\]
\deleted{\unsure{Supposons}} Alors $\delta$ est un groupe \uline{fini}; \inserted{\unsure{torsion de} $\operatorname{NS}(X_0)=\delta'$, soit $\delta''=\delta/\delta'$} et pour tout entier $n$ premier à $p$, on a une suite exacte canonique
\diagram{\begin{tikzcd}[column sep=tiny,row sep=small]
(*)\quad\unsure{0}\arrow[r,dash]&H^1(X,\mu_n)\arrow[r,"u_n"]\arrow[d,equal]&H^1(X_{\bar\eta},\mu_n)^I\arrow[r]&\deleted{\unsure{{}_n\delta''}}\inserted{{}_n\delta/{}_n\delta'}\arrow[r]&\deleted{\unsure{\delta'}}\delta'_n\arrow[r]&\delta_n\arrow[r]&\delta''_n\arrow[r]&0\\
&\inserted{H^1(X_0,\mu_n)}&&\deleted{\unsure{\delta'}}&&&&
\end{tikzcd}}
compatible avec les morphismes de transition évidents en $n$, d'où par passage à la limite, pour tout $\ell$ nbr pr. $\neq p$
\[
H^1(X,\mathbb Z_\ell)\xrightarrow{\sim}H^1(X_{\bar\eta},\mathbb Z_\ell)^I
\]
(\uline{Théorème des cycles \unsure{invariants|invariant}}).
\margin{(\unsure{cas $X_i$})}

\uline{Dém} $X$ est plat sur $S$, donc $X_\eta$ dense ds $X$, et $H^1(X,\mu_n)\hookrightarrow H^1(X_\eta,\mu_n)$. On trouve le diagramme
\diagram{\begin{tikzcd}[column sep=small,row sep=large]
H^1(X,\mu_n)\arrow[r,"\alpha","\sim"']\arrow[d]&{}_n\operatorname{Pic}(X)\arrow[d]&&\\
H^1(X_\eta,\mu_n)/(\mathbb Z/n\mathbb Z)\arrow[r,"\alpha_\eta","\sim"']\arrow[d]&{}_n\operatorname{Pic}(X_\eta)\arrow[r,"\beta_\eta","\sim"']\arrow[d]&{}_nP(K)\arrow[r,phantom,"\simeq"]\arrow[d]&{}_nP(\bar K)^I\arrow[dl,bend left=20]\\
H^1(X_{\bar\eta},\mu_n)\arrow[r,"\alpha_{\bar\eta}","\sim"']&{}_n\operatorname{Pic}(X_{\bar\eta})\arrow[r,"\beta_{\bar\eta}","\sim"']&{}_nP(\bar K)&
\end{tikzcd}}
où $\alpha,\alpha_\eta,\alpha_{\bar\eta}$ sont des isom par \deleted{\unsure{suite exacte}} de Kummer et $f_*(\underline O_X)\simeq\underline O_S$ (\uline{NB} $V^*/V^{*n}=0$, $K^*/K^{*n}\simeq\mathbb Z/n\mathbb Z$, car $V$ \uline{str.} local) $\beta_\eta$ (et $\beta_{\bar\eta}$) \unsure{sont} isom à cause de l'existence d'un 0-cycle de degré 1 sur $X_\eta$, donc on voit que $u_n$ dans (*) s'identifie:
\[
(**)\qquad{}_nv:\quad{}_n\operatorname{Pic}(X)\mathbin{-}{}_n\operatorname{Pic}(X_\eta)
\]
Comme $X$ est régulier (il suffirait: anneaux locaux factoriels) on a
\diagram{\begin{tikzcd}
0\arrow[r]&\Gamma\arrow[r,"i"]&\operatorname{Pic}(X)\arrow[r,dash,"v"]&\operatorname{Pic}(X_\eta)\arrow[r]&0
\end{tikzcd}}
où $\Gamma=\Delta/\mathbb Z\cdot X_0$, d'où
\[
\operatorname{Coker}u_n\simeq\operatorname{Coker}{}_nv\simeq\operatorname{Ker}i_n,\quad\text{où }i_n:\Gamma_n\longrightarrow\operatorname{Pic}(X)_n
\]

\mspage{64}
